\documentclass[11pt]{article}
\usepackage[T1]{fontenc}
\usepackage[margin=0.8in]{geometry}
\usepackage{lmodern,parskip,booktabs,tabularx,array,amsmath}
\usepackage{xcolor,hyperref,fancyhdr}
\hypersetup{colorlinks=true,urlcolor=blue,linkcolor=black}
\setlength{\headheight}{15pt}
\pagestyle{fancy}\fancyhf{}
\fancyhead[L]{Wenshan Zhang (wz3191)}\fancyhead[R]{DS-GA 1011: Assignment 4}
\fancyfoot[C]{\thepage}
\setcounter{secnumdepth}{0}
\newcolumntype{Y}{>{\raggedright\arraybackslash}X}
\begin{document}
\begin{center}{\Large\bfseries Assignment 4: Fine-tuning Language Models}\\[5pt]
Wenshan Zhang \quad NetID: wz3191
\end{center}
\textbf{Code:} \url{https://github.com/wz3191/dsga1011-assignment4}\\
\textbf{Final model checkpoint:} \href{https://drive.google.com/drive/folders/1FjvmzQhq-PaMB6C09AFjqRof1zkbUR1v}{Google Drive: final checkpoint and tokenizer}

\section{Q2a. Data transformation}
I introduce keyboard typos into the reviews. I treat each sequence of English letters as a word. Each word with at least four letters has a 35\% chance of being modified. For each selected word, I randomly choose one letter and replace it with its left or right neighbor in the same QWERTY row. Both the letter position and its replacement are chosen uniformly. At the end of a row, only one neighbor is available. I preserve capitalization, spaces, punctuation, and the original label.

This transformation is reasonable because users may press a nearby key when typing. Each selected word changes by only one letter, so readers can usually recognize the intended word. The overall sentiment should remain unchanged in most cases.

\section{Q3. Data augmentation}
I sampled 5,000 distinct training examples with seed 42, applied the same typo transformation, and added them to the original 25,000 examples. The augmented training set contained 30,000 examples.

\begin{center}
\begin{tabular}{lrr}\toprule
Training data & Original test accuracy & Typo test accuracy\\\midrule
Original & 93.072\% & 88.332\%\\
Original + augmented & 92.736\% & 89.912\%\\\bottomrule
\end{tabular}
\end{center}
Data augmentation improved typo test accuracy by 1.580 percentage points. Original test accuracy decreased by 0.336 percentage points. Training with similar typos may help the model recognize sentiment despite spelling errors. The added noise may slightly hurt performance on clean text. The small decrease could also reflect training variation.

One limitation is that training and evaluation use the same type of typo. These results do not show whether the model can handle other distribution shifts, such as slang or reviews from another domain.

\newpage
\section{Q4. Data statistics and processing}
I did not apply additional text preprocessing to the final model's data. Quotes and symbols have syntactic roles in SQL. Changing the case of string values may also change query results. The before and after statistics are therefore identical.

Statistics use the original \texttt{google-t5/t5-small} tokenizer, without padding, truncation, or special tokens. Vocabulary size counts distinct token IDs in each split, separately for natural language and SQL.

\begin{table}[h]\centering
\begin{tabular}{lrr}\toprule
Statistic & Train & Dev\\\midrule
Number of examples & 4225 & 466\\
Mean sentence length & 17.097 & 17.071\\
Mean SQL query length & 216.373 & 210.054\\
Vocabulary size (natural language) & 791 & 465\\
Vocabulary size (SQL) & 555 & 395\\\bottomrule
\end{tabular}
\caption{Data statistics before preprocessing.}
\end{table}
\begin{table}[h]\centering
\begin{tabular}{lrr}\toprule
Statistic & Train & Dev\\\midrule
Number of examples & 4225 & 466\\
Mean sentence length & 17.097 & 17.071\\
Mean SQL query length & 216.373 & 210.054\\
Vocabulary size (natural language) & 791 & 465\\
Vocabulary size (SQL) & 555 & 395\\\bottomrule
\end{tabular}
\caption{Data statistics after preprocessing for the final model.}
\end{table}

I also tested a separate tokenizer variant that added \texttt{<} and \texttt{<=} as ordinary tokens. Its mean SQL lengths were 216.207 (train) and 209.891 (dev), and its SQL vocabulary sizes were 556 and 396. The natural-language statistics did not change. This variant was not used in the final model.

\newpage
\section{Q5. T5 fine-tuning details}
\begin{table}[h]\centering\small
\renewcommand{\arraystretch}{1.15}
\begin{tabularx}{\textwidth}{p{0.19\textwidth}Y}\toprule
Design choice & Description\\\midrule
Data processing & Read one question and SQL query per line. Preserve text, case, punctuation, and example order. Shuffle training batches only. Do not add a task prefix. Test data contains questions only.\\
Tokenization & Use the original \texttt{google-t5/t5-small} tokenizer for questions and SQL. Include EOS (ID 1). Do not truncate. Dynamically pad each batch using PAD (ID 0), and mask encoder padding. Use \texttt{<extra\_id\_0>} (ID 32099) as the decoder start token. Decoder inputs contain this start token followed by the SQL tokens shifted right; targets contain SQL tokens followed by EOS.\\
Architecture & Initialize from pretrained \texttt{google-t5/t5-small}. Fine-tune the entire model, including shared embeddings, encoder, and decoder. No layers are frozen. The model has six encoder layers, six decoder layers, hidden size 512, and eight attention heads.\\
Training objective & Use teacher forcing and token cross-entropy. Exclude padded targets from the loss. Update the model after each batch.\\
Hyperparameters & AdamW; learning rate $10^{-4}$; weight decay 0; $\beta_1=0.9$, $\beta_2=0.999$; $\epsilon=10^{-8}$. Training and evaluation batch size 8. Linear learning-rate decay with no warmup. Seed 42. Train for at most 20 epochs; stop after five epochs without an improvement in dev Record F1.\\
Model selection & During training, evaluate with greedy decoding and save the checkpoint with the best dev Record F1. The selected checkpoint is epoch 17 (zero-based). Training reached epoch 19. Compare decoding methods on this same checkpoint.\\
Final decoding & Beam search with four beams, no sampling, and at most 512 new tokens. Use the same decoder start token as in training. Remove special tokens when decoding, disable tokenization-space cleanup, and strip leading and trailing whitespace.\\\bottomrule
\end{tabularx}
\caption{Details of the final T5 configuration.}
\end{table}

The final checkpoint is from the learning-rate experiment. It was trained from the pretrained T5 checkpoint, not continued from the earlier fine-tuned baseline. Beam search changes inference only; it does not update the model weights.

\newpage
\section{Q6. Results and error analysis}
\subsection{Quantitative results}
All development scores below are percentages. Execution errors include timeouts.
\begin{table}[h]\centering\small
\begin{tabular}{lrrrr}\toprule
Configuration & SQL EM & Record EM & Record F1 & Exec. errors\\\midrule
Baseline: $5\!\times\!10^{-5}$, greedy & 37.124 & 51.931 & 60.713 & 89/466\\
Tokenizer variant, greedy & 38.412 & 52.146 & 60.678 & 83/466\\
Baseline + beam 4 & 38.412 & 53.219 & 62.552 & 56/466\\
LR $10^{-4}$, greedy & 54.936 & 69.313 & 74.461 & 58/466\\
Full model: LR $10^{-4}$, beam 4 & 54.721 & 70.172 & 75.662 & 51/466\\\midrule
Full model: test & Pending & Pending & Pending & 60/432\\\bottomrule
\end{tabular}
\caption{Development results and test evaluation status.}
\end{table}

\textbf{Learning-rate ablation.} Compared with the original baseline, the learning-rate experiment changed the initial learning rate from $5\times10^{-5}$ to $10^{-4}$. Other training settings and the original tokenizer stayed the same. With greedy decoding, dev Record F1 increased by 13.748 percentage points.

\textbf{Decoding ablation.} With the baseline checkpoint fixed, beam 4 improved Record F1 by 1.839 points over greedy decoding. With the LR $10^{-4}$ checkpoint fixed, beam 4 improved Record F1 by 1.200 points. These comparisons change only decoding. The combined configuration had the highest dev Record F1, so I selected it for test generation.

\textbf{Tokenizer experiment.} Adding the SQL operators and adjusting their decoded spacing did not improve Record F1 over the original baseline. I therefore kept the original tokenizer. This experiment changed the tokenizer and its associated decoding handling together. All comparisons use one seed, so they do not establish statistical significance.

Record EM improved with the final beam setting, while SQL EM decreased slightly. Different SQL strings can return the same records. Also, beam search optimizes model sequence scores, not Record F1 directly.

Test scores are pending Gradescope evaluation.

\newpage
\subsection{Qualitative error analysis}
I analyzed the full model's 466 dev predictions. Line numbers below are one-based in \texttt{data/dev.nl}. I first classified execution errors by their saved error message. Among queries without execution errors, I then compared predicted and reference record sets. The four categories below are disjoint.

\begin{table}[h]\centering\small
\renewcommand{\arraystretch}{1.2}
\begin{tabularx}{\textwidth}{p{0.15\textwidth}Y Y p{0.105\textwidth}}\toprule
Error type & Example & Description and possible reason & Count/total\\\midrule
Syntax or incomplete SQL & Line 25 asks for arrival before 10am. The prediction contains \texttt{flight\_1.arrival\_time 1000}, without an operator. & SQLite reports a syntax error. The original tokenizer cannot preserve every SQL symbol, which can contribute to missing operators. Other predictions have extra parentheses or incomplete SQL. & 34/466\\
Invalid or ambiguous columns & Line 98 asks about Pittsburgh ground transportation. The prediction uses \texttt{city\_2.city\_code}, but its FROM clause does not declare \texttt{city\_2}. & The model does not consistently connect column references with declared table aliases. Unconstrained generation does not enforce schema validity. & 16/466\\
Incorrect returned records & Line 6: ``I'm starting from Denver.'' The prediction requires both the origin and destination to be Denver. It returns 0 records; the reference returns 611. & The query runs, but adds an unsupported condition. This category also includes other executable queries with incorrect result sets. Common query patterns may lead the model to add or miss conditions. & 96/466\\
Execution timeout & Line 24 asks for flights and prices from LA to Charlotte. The prediction lists fare tables but omits their joins to the flight table. & The saved result reports a timeout. Missing join conditions can create large Cartesian products, making execution expensive. & 1/466\\\bottomrule
\end{tabularx}
\caption{Error analysis for the final LR $10^{-4}$ + beam 4 model.}
\end{table}

There are 319 other queries that execute without errors and return the exact reference record set. These category counts are diagnostic counts, not a replacement for the official metrics. In particular, the official evaluator can count an empty result from an execution error as an exact record match when the reference result is also empty.

\end{document}
